\begin{table}[!tbp]
\centering
\caption{Paired annual physical contrasts at $\eta=1$. Effects are candidate minus comparator in the unit shown; intervals and the LOO range are computed as in Table~\ref{tab:r2-contrasts}.}
\label{tab:r2-physical-contrasts}
\scriptsize
\setlength{\tabcolsep}{3pt}
\begin{tabular}{@{}lllrrr@{}}
\toprule
Contrast & Endpoint & Unit & Effect & 95\% interval & LOO range \\
\midrule
P $-$ none & Uncovered & MWh/yr & -338.43 & [-868.05, -23.29] & [-394.83, -79.54] \\
P $-$ none & $W_7$ deficit & MWh & -89.60 & [-246.07, 0.00] & [-104.53, -13.25] \\
P $-$ none & $W_7$ adequacy & fraction & 0.0043 & [0.0000, 0.0094] & [0.0020, 0.0050] \\
P $-$ none & Extra booked & $10^6$ MW$\cdot$h/yr & 1.1236 & [0.9446, 1.2924] & [1.0745, 1.1855] \\
P $-$ F2 & $W_7$ deficit & MWh & 193.51 & [12.03, 430.36] & [92.67, 225.76] \\
P $-$ F2 & Uncovered & MWh/yr & 105.35 & [-253.56, 444.58] & [-10.19, 260.99] \\
\bottomrule
\end{tabular}
\par\footnotesize A negative uncovered or deficit effect is better; a positive adequacy effect is better. The $W_7$ deficit and adequacy intervals touch zero on their favorable side, while the annual uncovered reduction excludes zero; the $(P-F2)$ rows keep the same orientation and show that the production rule leaves a larger $W_7$ deficit than bank pacing. Inference follows Section~\ref{sec:protocol}.
\end{table}
